\scanpage{6}

\mspage{4}

\uline{Proposition 2}. Soit $\Phi$ une partie close à g. de $\mathrm{Fl}\,M$. Alors $\Phi$ satisfait les conditions a) b) c) ci-dessous :

a) $\Phi$ contient les iso, et est stable par composition, et par cochangement de base (i.e si on a

\diagram{\begin{tikzcd}[ampersand replacement=\&,execute at end picture={\draw[line width=.4pt] (oldx.south west) -- (oldxp.north west) -- (oldy.west) -- (oldxp.east) -- (oldy.north east) -- (oldyp.north east) -- (oldyp.south west);}]
|[alias=oldx]|\deleted{X}\arrow[r] \& |[alias=oldxp]|\deleted{X'}\\
|[alias=oldy]|\deleted{Y}\arrow[u,"\deleted{\unsure{f|j}}"]\arrow[r,"\deleted{i}"'] \& |[alias=oldyp]|\deleted{\unsure{Y'|X'}}\arrow[u,"\deleted{\unsure{f'|j'}}"']
\end{tikzcd}}

\diagram{\begin{tikzcd}[ampersand replacement=\&]
A\arrow[r,"i"]\arrow[d] \& B\arrow[d]\\ A'\arrow[r,"i'",no head]\& B'
\end{tikzcd}}

carré cocartésien \unsure{dans|de} $M$, avec $i\in\Phi$, on a $i'\in\Phi$).

b) $\Phi$ stable par \unsure{facteur direct|facteurs directs} (en tant que partie de l'ens des objets de $\underline{\mathrm{Fl}}(M)$).

c) $\Phi$ \deleted{\unsure{est} $\unsure{\Phi}$} ``stable par limites inductives ordinales'', dans le sens suivant : Soit $\deleted{\unsure{A}}I$ un ens bien ordonné, $(A_i)_{i\in I}$ un \deleted{système} inductif dans $M$, tel que les $u_{\unsure{ji|ij}}:A_i\unsure{-|\to}A_j$ soient dans $\Phi$, et $A_{i_0}\xleftarrow{\sim}\varinjlim_{i<i_0} A_i$ si $i_0\in I$ est un élément limite. Supposons que $A=\varinjlim_I A_i$ existe. Alors $\forall i_0\in I$, la flèche canonique $A_{\unsure{i_0|j_0}}\longrightarrow A$ est dans $\Phi$.

Il y a un énoncé dual pour \unsure{une|les} partie\deleted{\unsure{s}} $\Psi$ close\deleted{\unsure{s}} à droite\unsure{;|'';} \deleted{\unsure{cont.}} elle contient les \unsure{izo|iso}, est stable par composition et par chgt de base, \inserted{par facteurs directs, \unsure{et enfin}} et stable par limites projectives ordinales.

% archive sheet number: 5

